\chapter{\MakeUppercase{Discussão}}

\section{Acesso a dados e transparência pública}

O SIM materializa o princípio de que dados abertos só cumprem sua função democrática quando acessíveis \cite{Brasil2016DadosAbertos,Lei2012LAI}. Ao integrar consulta analítica e recuperação documental em canal conversacional amplamente utilizado no Brasil, a plataforma reduz a distância entre a publicação formal de bases governamentais e o uso efetivo por cidadãos, jornalistas e organizações da sociedade civil.

\section{Comparação com benchmarks acadêmicos}

A equivalência funcional de 72,5\% (Prompt B) e 62,5\% (E1) \textbf{não é diretamente comparável} à EX de Spider ou BIRD, pelas razões seguintes:

\begin{itemize}
    \item \textbf{Métrica distinta}: Spider/BIRD reportam EX estrita (multiconjunto de pares coluna--valor); este trabalho adota EF com mapeamento de colunas de resposta e tolerância $\tau = 0{,}5\%$ em valores fiscais. No E1, EX estrita foi de apenas 10\%, enquanto EF foi de 62,5\% --- evidência de que EX estrita subestimaria respostas úteis ao cidadão;
    \item \textbf{Domínio e escala}: Spider reúne 200 bases em inglês com esquemas sintéticos; BIRD incorpora bases reais volumosas com ruído. O SIM opera sobre esquema SICOM único, em português, com regras contábeis locais --- tarefa mais estreita em diversidade de esquema, porém mais fiel ao uso municipal;
    \item \textbf{Ordem de grandeza}: em BIRD, GPT-4 reporta EX na faixa de 50--55\% \cite{Li2024BIRD}; métodos decomposicionais (DIN-SQL) elevam para $\approx$60\% \cite{Pourreza2024DIN}. A EF de 72,5\% do SIM situa-se na mesma ordem de magnitude, com ressalva de que EF é métrica mais permissiva que EX;
    \item \textbf{Amostra}: o conjunto de referência SQL tem $n=80$ casos correlacionados (10 intenções $\times$ variações), não 80 perguntas independentes --- intervalos de confiança amplos (IC EF do Prompt B: 61,9--81,1\%) refletem incerteza amostral.
\end{itemize}

Em síntese, o SIM não compete em \textit{leaderboard} acadêmico, mas demonstra viabilidade funcional em dados municipais reais com latência compatível a WhatsApp --- objetivo distinto dos \textit{benchmarks} generalistas.

\section{Achados empíricos principais}

Três achados estruturam a discussão, sem repetir numericamente o Capítulo~4:

\begin{enumerate}
    \item \textbf{Engenharia de contexto supera troca de modelo}: Prompt B (+10 pp EF) superou GPT-4o (E4) em EF com fração da latência --- reforçando que regras de domínio SICOM são alavanca principal neste cenário;
    \item \textbf{Recuperação híbrida é necessária para atos numerados}: no estrato \texttt{com\_numero}, busca densa isolada falhou (3,1\% Recall@5); lookup por número/ano no cabeçalho do \textit{chunk} resolve o caso de uso mais frequente em Diário Oficial;
    \item \textbf{Falhas residuais concentram-se em cruzamentos}: EF em consultas ``difíceis'' (57,1\%) e cruzamentos investigativos (30\%) indica que \textit{joins} entre SICOM e CGU por nome ou documento permanecem desafio mesmo com Prompt B.
\end{enumerate}

\section{Configuração adotada na instância Caeté}

Com base nos três lotes experimentais, a instância Caeté adota: (i)~GPT-4.1-mini com Prompt B para Text-to-SQL; (ii)~recuperação híbrida com deduplicação e filtro por tipo de ato para o subsistema RAG. Modelos alternativos e \textit{prompt} estendido foram descartados por desempenho inferior ou latência incompatível com canal conversacional síncrono.

\section{Confiabilidade, risco e segurança}

Mesmo na configuração final (EF 72,5\%), aproximadamente um quarto das consultas do conjunto de referência não produz grandezas funcionalmente equivalentes ao gabarito. Em canal público voltado ao cidadão, isso implica risco de respostas fiscalmente incorretas --- especialmente em cruzamentos investigativos e índices derivados. Mitigações recomendadas incluem: (i)~exibir a consulta SQL gerada ou um resumo auditável quando a pergunta envolve valores monetários; (ii)~mensagens de incerteza quando a confiança do resultado é baixa; e (iii)~encaminhamento a fontes oficiais (portal de transparência, Diário Oficial).

Quanto à segurança, um canal conversacional público está exposto a \textit{prompt injection} e, indiretamente, a SQL gerada sobre base local. O uso de API com permissões restritas, validação sintática antes da execução e ausência de comandos de escrita no DuckDB reduzem, mas não eliminam, o risco. A camada de resposta em linguagem natural pode ainda alucinar interpretações sobre resultados corretos --- reforçando a necessidade de transparência do dado tabular subjacente.

\section{Limitações}
\label{sec:limitacoes}

\begin{itemize}
    \item O conjunto SQL contém 80 casos derivados de 10 intenções analíticas com variações paramétricas correlacionadas;
    \item O ganho do Prompt B foi medido no mesmo conjunto usado para diagnosticar falhas; as regras são portáveis ao padrão SICOM, mas a magnitude do ganho pode variar em outras instâncias (McNemar: $p = 0{,}146$);
    \item Conjunto de referência de uma única instância (Caeté);
    \item Lacuna de folha de pagamento 2020--2022 na base local;
    \item Gabarito manual pode não cobrir todas as SQL semanticamente corretas;
    \item Definição automatizada de colunas de resposta em 75 dos 80 itens SQL;
    \item Julgamento automático (JAR) não discriminante na amostra piloto ($n=30$);
    \item Demonstração WhatsApp sem métricas quantitativas de adoção em produção.
\end{itemize}

\section{Implicações práticas}

O SIM demonstra viabilidade técnica de consulta conversacional a dados municipais abertos, com arquitetura replicável para outras prefeituras que operem sobre o padrão SICOM. Os resultados empíricos subsidiam a escolha de modelo, engenharia de contexto e estratégia de recuperação documental em instâncias municipais com perfil de dados semelhante ao de Caeté.
